\documentclass[onecolumn]{IEEEtran}

% -------------------------------------------------
% Packages
% -------------------------------------------------

\usepackage{amsmath}
\usepackage{amssymb}
\usepackage{algorithm}
\usepackage{algorithmic}
\usepackage{array}
\usepackage{booktabs}
\usepackage{cite}
\usepackage{url}
\usepackage[T1]{fontenc}
\usepackage{graphicx}
\usepackage{epstopdf}

% -------------------------------------------------
% Document
% -------------------------------------------------

\begin{document}

% -------------------------------------------------
% Title
% -------------------------------------------------

\title{
RICAN: A Fixed-Rate Binary Representation Framework for
Finite Alphabets Using Continued-Fraction Convergents
}

\author{
\IEEEauthorblockN{
J. Jes\'us Mart\'inez Palomo
}
\IEEEauthorblockA{
Independent Researcher\\
Email: enviarcorreoajesusdelaunam@gmail.com
}
}

\maketitle

% -------------------------------------------------
% Abstract
% -------------------------------------------------

\begin{abstract}

This paper introduces RICAN (Representational Isomorphic Coding for
Arbitrary Numeration), a mathematical framework for constructing
asymptotically tight fixed-rate reversible binary representations for
sequences over arbitrary finite alphabets.

The central contribution is a constructive method for selecting block
sizes through continued-fraction approximation of $\log_2(B)$.
The method identifies convergent-derived block lengths that generate
deterministic fixed-rate representations with arbitrarily small absolute
rounding gaps along suitable subsequences. Consequently, the normalized
per-symbol representation overhead converges to zero, where $B$ denotes
the alphabet cardinality.

RICAN does not perform source compression; it minimizes deterministic
fixed-rate binary representation overhead.

For a block of length $K$ over an alphabet of size $B$, the minimum
number of bits required by any fixed-length injective binary
representation is bounded by

\[
n(K)=\lceil K\log_2(B)\rceil
\]

bits.

The unavoidable rounding overhead is therefore

\[
G_B(K)=
\lceil K\log_2(B)\rceil-K\log_2(B).
\]

RICAN demonstrates that denominators of continued-fraction
convergents of $\log_2(B)$ provide an infinite subsequence of block
sizes whose absolute rounding gap approaches zero, and therefore
whose normalized per-symbol representation overhead converges to
zero.

The framework applies uniformly to the infinite family of finite
alphabets

\[
\mathcal{A}=\{\Sigma_B:B\geq2\},
\]

where

\[
\Sigma_B=\{0,1,\ldots,B-1\}.
\]

Power-of-two alphabets achieve exact binary representations with zero
overhead. For non-power-of-two alphabets, continued-fraction
convergents provide deterministic asymptotically tight block sizes.

For the decimal alphabet ($B=10$), the practical reference
implementation uses the convergent-derived block size

\[
K=643,
\]

requiring

\[
n(K)=2136
\]

bits and achieving a representation overhead of approximately

\[
3.65\times10^{-7}
\]

bits per symbol.

A larger upper convergent denominator

\[
K=783032
\]

corresponds to the convergent

\[
\frac{2601184}{783032}
\]

in the continued fraction expansion of $\log_2(10)$, satisfying

\[
K\log_2(10)=2601183.99999948\ldots,
\]

requiring

\[
n(K)=\lceil K\log_2(10)\rceil=2601184
\]

bits, represented exactly by

\[
325148
\]

bytes.

The total rounding gap of this complete block is approximately
$5.20\times10^{-7}$ bits, while the normalized per-symbol overhead is

\[
\epsilon_{10}(783032)
=
\frac{5.20\times10^{-7}}{783032}
\approx
6.64\times10^{-13}
\]

bits per symbol.

In this context, representation proximity refers exclusively to the
binary realization overhead and does not imply source compression.

Unlike statistical source coding methods, RICAN does not exploit
source redundancy or probability distributions. The encoded length
depends only on alphabet cardinality and input length.

The main contribution is a constructive procedure for selecting
fixed-rate block representations with asymptotically vanishing
absolute rounding gap through continued-fraction approximation.

A reference implementation using the RJ03 self-contained
representation format is provided.

\end{abstract}

% -------------------------------------------------
% Index Terms
% -------------------------------------------------

\begin{IEEEkeywords}

Fixed-rate representation,
finite alphabets,
continued fractions,
Diophantine approximation,
enumerative representation,
bijective coding.

\end{IEEEkeywords}

% =================================================
% I. INTRODUCTION
% =================================================

\section{Introduction}

Lossless information systems require reversible mappings between
symbol sequences and binary descriptions.

Two different objectives are frequently associated with such mappings:
source compression and binary representation.

Compression methods attempt to reduce expected description length by
exploiting statistical properties of the source. In contrast,
representation methods study the deterministic mapping between a finite
alphabet and a binary storage space.

RICAN belongs to the second category.

The objective of RICAN is not to compress information content, but to
minimize the unavoidable integer rounding overhead that appears when
finite alphabet spaces are represented using binary hardware.

For an alphabet containing $B$ symbols, the information quantity
associated with one uniformly selected symbol is

\[
\log_2(B)
\]

bits.

When

\[
B=2^m ,
\]

the representation is exact because

\[
\log_2(B)=m .
\]

However, for alphabet cardinalities that are not powers of two,
$\log_2(B)$ is generally irrational, preventing an exact
one-symbol fixed-rate binary representation.

A block representation solves this problem by grouping symbols.

For a block length $K$, the number of possible sequences is

\[
|\Sigma_B^K|=B^K .
\]

A binary representation containing $n$ bits provides

\[
2^n
\]

possible states.

Therefore, any reversible fixed-rate representation must satisfy

\[
2^n\geq B^K .
\]

Taking logarithms gives:

\[
n\geq K\log_2(B).
\]

Because the number of bits must be an integer, the minimum required
block length is:

\[
n(K)=\lceil K\log_2(B)\rceil .
\]

The resulting rounding overhead is:

\[
G_B(K)=
\lceil K\log_2(B)\rceil
-
K\log_2(B).
\]

The central problem addressed by RICAN is the selection of block sizes
$K$ that minimize this deterministic gap.

The contribution of this work is a constructive method based on
continued fractions.

For irrational values of $\log_2(B)$, the denominators of
continued-fraction convergents provide block lengths where the product

\[
K\log_2(B)
\]

approaches an integer with increasing accuracy.

The framework is defined over the infinite family of finite alphabets:

\[
\mathcal{A}
=
\{\Sigma_B:B\geq2\}.
\]

Therefore, binary alphabets, DNA alphabets, decimal digits,
hexadecimal symbols, byte alphabets, and arbitrary finite alphabets
are treated as particular instances of the same mathematical model.

The main contributions of this paper are:

\begin{enumerate}

\item A general fixed-rate binary representation mathematical framework
for finite alphabets.

\item A constructive block-size selection procedure based on
continued-fraction convergents of $\log_2(B)$.

\item A theoretical analysis proving the existence of convergent
subsequences with vanishing absolute representation gap.

\item A self-contained reversible binary representation format (RJ03)
that stores all information required for deterministic reconstruction.

\item An experimental validation over decimal sequences and multiple
finite alphabets.

\end{enumerate}

The remainder of this paper is organized as follows.

Section II reviews related approaches.

Section III introduces the mathematical foundations.

Section IV presents the main theoretical results.

Section V describes the block selection algorithm.

Section VI presents the RICAN representation framework and RJ03 format.

Section VII presents experimental evaluation.

Section VIII discusses implications and limitations.

Section IX presents code availability.

Section X concludes the paper.

% =================================================
% NOTATION
% =================================================

\section{Notation}

The following notation is used throughout this paper.

\begin{table}[!t]
\caption{Notation Summary}
\label{tab:notation}
\centering
\begin{tabular}{ll}
\toprule
Symbol & Meaning\\
\midrule
$B$ & Alphabet cardinality\\
$\Sigma_B$ & Alphabet of size $B$\\
$K$ & Block length in symbols\\
$n(K)$ & Binary representation length of a block\\
$G_B(K)$ & Absolute representation gap\\
$\delta_B(K)$ & Normalized representation overhead\\
$\epsilon_B(K)$ & Per-symbol representation overhead\\
$p_n/q_n$ & Continued-fraction convergent\\
$L$ & Input sequence length\\
$N(L)$ & Total representation length\\
\bottomrule
\end{tabular}
\end{table}

The notation emphasizes the distinction between source size and
representation size.

The alphabet cardinality $B$ determines the ideal rate
$\log_2(B)$, while the block length $K$ determines the integer
rounding behavior.

% =================================================
% II. RELATED WORK
% =================================================

\section{Related Work}

The problem of representing finite-alphabet information using binary
hardware has been studied from several perspectives, including
statistical coding, enumerative representation, and fixed-length
mapping techniques.

RICAN differs from these approaches because its objective is not the
reduction of expected source description length, but the minimization
of deterministic binary representation overhead.

\subsection{Statistical Source Coding}

Classical source coding methods optimize representation length when
statistical properties of the source are available.

Huffman coding assigns variable-length prefix codewords according to
symbol probabilities and provides optimality among binary prefix codes
for a known distribution \cite{huffman1952}.

Arithmetic coding represents sequences as subintervals of the unit
interval and avoids the integer-length limitation of conventional
prefix coding \cite{rissanen1979}.

Asymmetric numeral systems (ANS) achieve entropy coding performance
comparable to arithmetic coding while allowing efficient finite-state
implementations \cite{duda2013}.

These techniques exploit statistical redundancy.

Their objective is therefore different from RICAN, which constructs
fixed-rate representations whose size depends only on alphabet
cardinality and sequence length.

\subsection{Dictionary-Based Methods}

Dictionary-based methods reduce source description length by identifying
repeated structures in sequences.

The Lempel-Ziv family introduced universal algorithms that represent
repeated substrings through references to previously observed patterns
\cite{ziv1977}.

Modern practical compression systems frequently combine dictionary
methods with entropy coding to improve performance on structured data.

RICAN does not perform pattern detection, probability estimation, or
redundancy removal.

Instead, it addresses the independent problem of representing every
possible block of a finite alphabet using a deterministic binary
mapping with minimum integer rounding overhead.

\subsection{Enumerative Representation}

Enumerative methods provide systematic ranking and unranking procedures
for finite combinatorial spaces.

Cover introduced enumerative source encoding as a method for mapping
sets of sequences into integer indices
\cite{cover1973}.

Later developments extended these ideas to constrained sequences and
efficient combinatorial representations
\cite{immink2004}.

RICAN is related to enumerative representation because both approaches
use the cardinality of finite spaces.

However, the primary objective differs.

Enumerative coding focuses on ranking specific classes of sequences,
often under constraints.

RICAN focuses on the selection of block lengths $K$ such that the
binary representation length

\[
\lceil K\log_2(B)\rceil
\]

has minimum deterministic rounding error.

\subsection{Fixed-Length Binary Representation}

A finite alphabet with cardinality $B$ requires at least

\[
\lceil\log_2(B)\rceil
\]

bits when symbols are represented independently.

This approach is optimal only when $B$ is a power of two.

For general alphabets, block representation reduces the effect of
rounding by considering the complete block space:

\[
|\Sigma_B^K|=B^K .
\]

The minimum binary length is therefore:

\[
n(K)=\lceil K\log_2(B)\rceil .
\]

RICAN studies the mathematical selection of $K$ rather than the
construction of a new coding probability model.

\subsection{Continued Fractions and Diophantine Approximation}

Continued fractions provide highly accurate rational approximations to
irrational numbers.

For an irrational number $x$, the continued fraction expansion is

\[
x=[a_0;a_1,a_2,\ldots].
\]

The convergents are fractions

\[
\frac{p_n}{q_n}
\]

with the classical approximation property:

\[
\left|
x-\frac{p_n}{q_n}
\right|
<
\frac{1}{q_nq_{n+1}} .
\]

Multiplying by $q_n$ gives:

\[
|q_nx-p_n|
<
\frac{1}{q_{n+1}} .
\]

Therefore, the denominators $q_n$ identify integer multiples of
$x$ that approach integer values with increasing accuracy.

RICAN applies this classical number-theoretic property to the
selection of representation block sizes.

For

\[
x=\log_2(B),
\]

the selected denominators provide block lengths where:

\[
K\log_2(B)
\]

is close to an integer boundary.

\subsection{Position of RICAN}

RICAN should be interpreted as a mathematical representation
framework.

Its contribution is not an entropy coding algorithm and not a
compression mechanism.

The framework provides:

\begin{itemize}

\item a general model for finite alphabet representation,

\item a method for selecting fixed-rate block sizes,

\item a theoretical analysis based on continued-fraction approximation,

\item a self-contained reversible binary format.

\end{itemize}

The conceptual distinction between RICAN and conventional coding
methods is summarized in Table~\ref{tab:comparison}.

\begin{table}[!t]
\caption{Conceptual Comparison Between Representation and Compression
Approaches}
\label{tab:comparison}
\centering
\begin{tabular}{lll}
\toprule
Method & Objective & Statistics Required\\
\midrule
Huffman Coding &
Expected length minimization &
Yes\\

Arithmetic Coding &
Entropy-rate approximation &
Yes\\

ANS &
Entropy coding efficiency &
Yes\\

Dictionary Methods &
Pattern redundancy reduction &
Yes\\

Enumerative Coding &
Finite-space ranking &
Sometimes\\

RICAN &
Fixed-rate representation gap minimization &
No\\

\bottomrule
\end{tabular}
\end{table}

% =================================================
% III. PRELIMINARIES
% =================================================

\section{Preliminaries}

This section introduces the mathematical definitions required for the
RICAN framework.

\subsection{Finite Alphabet Family}

RICAN considers the family of finite alphabets:

\[
\mathcal{A}
=
\{\Sigma_B:B\geq2\},
\]

where:

\[
\Sigma_B=
\{0,1,\ldots,B-1\}.
\]

The value $B$ completely determines the representation problem.

Representative examples include:

\begin{itemize}

\item $B=2$: binary alphabet.

\item $B=4$: DNA alphabet.

\item $B=10$: decimal digits.

\item $B=16$: hexadecimal alphabet.

\item $B=256$: byte alphabet.

\end{itemize}

\subsection{Rationality of $\log_2(B)$}

\begin{lemma}
\label{lem:rationality}

For an integer $B\geq2$,

\[
\log_2(B)\in\mathbb{Q}
\]

if and only if

\[
B=2^m
\]

for some integer $m\geq1$.

\end{lemma}

\begin{proof}

If:

\[
B=2^m,
\]

then:

\[
\log_2(B)=m,
\]

which is rational.

Conversely, assume:

\[
\log_2(B)=\frac{p}{q}
\]

in lowest terms.

Then:

\[
B^q=2^p .
\]

By unique factorization, $B$ can contain only the prime factor $2$.
Therefore:

\[
B=2^m
\]

for some integer $m$.

\end{proof}

\subsection{Block Space}

For a block length $K$, the set of possible input blocks is:

\[
\Sigma_B^K .
\]

Its cardinality is:

\[
|\Sigma_B^K|=B^K .
\]

A binary word of length $n$ provides:

\[
2^n
\]

possible states.

Therefore, an injective representation requires:

\[
2^n\geq B^K .
\]

Taking logarithms:

\[
n\geq K\log_2(B).
\]

Because $n$ must be an integer:

\[
n(K)=\lceil K\log_2(B)\rceil .
\]

\subsection{Absolute Representation Gap}

\begin{definition}

For an alphabet cardinality $B$, the absolute representation gap is:

\[
G_B(K)=
\lceil K\log_2(B)\rceil
-
K\log_2(B).
\]

This quantity measures the integer rounding distance, in bits, between
the ideal real-valued representation length and the actual binary block
length.

\end{definition}

\subsection{Normalized Representation Overhead}

\begin{definition}

The normalized representation overhead is:

\[
\delta_B(K)
=
\frac{G_B(K)}
{K\log_2(B)} .
\]

The per-symbol representation overhead is:

\[
\epsilon_B(K)
=
\frac{G_B(K)}{K}.
\]

These quantities describe representation tightness relative to the
ideal fixed-rate bound.

They do not represent compression efficiency.

\end{definition}

\subsection{Asymptotically Tight Representation}

\begin{definition}

A sequence of block sizes $K_i$ defines an asymptotically tight
representation with respect to rounding overhead if:

\[
\lim_{i\rightarrow\infty}
\delta_B(K_i)=0 .
\]

For RICAN, the selected sequence is obtained from suitable
continued-fraction convergent denominators of $\log_2(B)$.

\end{definition}

\subsection{Fixed-Rate Representation Length}

For an input sequence of length $L$, the total representation length
is:

\[
N(L)=
\left\lceil\frac{L}{K}\right\rceil n(K),
\]

where:

\[
n(K)=\lceil K\log_2(B)\rceil .
\]

This expression assumes the final incomplete block is encoded using
its actual bit length and excludes constant container metadata.

For large $K$ and large $L$, the effective rate approaches:

\[
\log_2(B)
\]

bits per symbol along the selected convergent subsequence.

\subsection{Mathematical Relation to Enumerative Coding}

RICAN and enumerative coding share the use of integer ranking over
finite spaces.

In enumerative coding, sequences are mapped to integer indices through
combinatorial counting.

In RICAN, the mapping is defined by radix conversion over complete
blocks:

\[
X(S)=\sum_{i=0}^{K-1}s_iB^{K-1-i}.
\]

Both approaches produce reversible mappings.

However, the optimization criterion differs.

Enumerative coding typically minimizes expected index length under
probability constraints.

RICAN minimizes deterministic rounding overhead at the block level
through continued-fraction block selection.

Therefore, RICAN can be interpreted as a special case of enumerative
representation where:

\begin{itemize}

\item the enumerated space is the complete Cartesian product
$\Sigma_B^K$,

\item the objective is block-length optimization,

\item the selection of $K$ is performed by Diophantine approximation.

\end{itemize}

This connection clarifies the position of RICAN within existing
representation techniques.

% =================================================
% IV. MAIN THEORETICAL RESULTS
% =================================================

\section{Main Theoretical Results}

This section establishes the mathematical foundation of RICAN.

The central observation is that the fixed-rate representation overhead
depends on the distance between:

\[
K\log_2(B)
\]

and the nearest integer value.

For an alphabet cardinality $B$, define:

\[
x_B=\log_2(B).
\]

The representation gap is:

\[
G_B(K)
=
\lceil Kx_B\rceil-Kx_B .
\]

The objective of the RICAN construction is to select block lengths
$K$ for which this quantity becomes arbitrarily small.

\subsection{Distance to Integer Values}

For a real number $y$, define its distance to the closest integer as:

\[
\operatorname{dist}(y,\mathbb{Z})
=
\min_{m\in\mathbb{Z}}|y-m|.
\]

For:

\[
x_B=\log_2(B),
\]

define:

\[
d_B(K)
=
\operatorname{dist}(Kx_B,\mathbb{Z}).
\]

This quantity measures how close the ideal representation length is to
an integer number of bits.

When the closest integer corresponds to the ceiling operation:

\[
\lceil Kx_B\rceil,
\]

the representation gap satisfies:

\[
G_B(K)=d_B(K).
\]

\subsection{Continued-Fraction Approximation}

Let $x$ be irrational and let:

\[
\frac{p_n}{q_n}
\]

be a convergent of its continued-fraction expansion.

A classical property of continued fractions is:

\[
\left|
x-\frac{p_n}{q_n}
\right|
<
\frac{1}{q_nq_{n+1}} .
\]

Multiplying by $q_n$ gives:

\[
|q_nx-p_n|
<
\frac{1}{q_{n+1}} .
\]

Because $p_n$ is an integer:

\[
\operatorname{dist}(q_nx,\mathbb{Z})
\leq
|q_nx-p_n|.
\]

Therefore:

\[
\operatorname{dist}(q_nx,\mathbb{Z})
<
\frac{1}{q_{n+1}} .
\]

\begin{lemma}
\label{lem:convergent}

Let $x$ be irrational and let

\[
\frac{p_n}{q_n}
\]

be a continued-fraction convergent of $x$.

Then:

\[
\operatorname{dist}(q_nx,\mathbb{Z})
<
\frac{1}{q_{n+1}} .
\]

\end{lemma}

\begin{proof}

The classical approximation theorem for continued fractions gives:

\[
|q_nx-p_n|
<
\frac{1}{q_{n+1}} .
\]

Since $p_n$ belongs to $\mathbb{Z}$,

\[
\operatorname{dist}(q_nx,\mathbb{Z})
\leq
|q_nx-p_n|.
\]

Combining both expressions proves the result.

\end{proof}

\subsection{Convergent-Based Block Selection}

The previous lemma directly provides a construction rule for RICAN.

For:

\[
x_B=\log_2(B),
\]

choose block sizes:

\[
K=q_n,
\]

where $q_n$ are denominators of continued-fraction convergents.

Then:

\[
Kx_B
\]

approaches an integer value with an error bounded by:

\[
\frac{1}{q_{n+1}} .
\]

\begin{theorem}
\label{thm:convergent}

Let:

\[
x_B=\log_2(B)
\]

be irrational.

Let:

\[
\frac{p_n}{q_n}
\]

be the convergents of the continued-fraction expansion of $x_B$.

Then the block lengths:

\[
K_n=q_n
\]

satisfy:

\[
\operatorname{dist}(K_nx_B,\mathbb{Z})
<
\frac{1}{q_{n+1}} .
\]

Consequently, there exists an infinite sequence of RICAN block sizes
whose ideal representation length approaches an integer value with
arbitrarily small absolute error.

\end{theorem}

\begin{proof}

From Lemma~\ref{lem:convergent}:

\[
\operatorname{dist}(q_nx_B,\mathbb{Z})
<
\frac{1}{q_{n+1}} .
\]

The denominators of continued-fraction convergents satisfy:

\[
q_n\rightarrow\infty .
\]

Therefore:

\[
\frac{1}{q_{n+1}}\rightarrow0 .
\]

Hence:

\[
\lim_{n\rightarrow\infty}
\operatorname{dist}(q_nx_B,\mathbb{Z})
=0 .
\]

This proves that the convergent denominators provide an infinite
sequence of block sizes with vanishing distance between the ideal and
integer representation lengths.

\end{proof}

\subsection{Existence of Vanishing Absolute Representation Gap}

The previous theorem concerns distance to the nearest integer.

For RICAN, the relevant quantity is the ceiling gap:

\[
G_B(K)
=
\lceil Kx_B\rceil-Kx_B .
\]

\begin{theorem}
\label{thm:abs_overhead}

Let:

\[
x_B=\log_2(B)
\]

be irrational.

Then there exists an infinite subsequence of continued-fraction
convergent denominators:

\[
K_i=q_{n_i}
\]

such that:

\[
G_B(K_i)\rightarrow0 .
\]

Therefore, the absolute fixed-rate representation overhead of RICAN
can be made arbitrarily small through convergent-based block selection.

\end{theorem}

\begin{proof}

From Theorem~\ref{thm:convergent}, the convergent denominators satisfy:

\[
\operatorname{dist}(q_nx_B,\mathbb{Z})
<
\frac{1}{q_{n+1}} .
\]

The convergents alternate around the irrational number $x_B$.
Therefore, infinitely many convergents approach $x_B$ from above and
infinitely many from below.

Consider the subsequence approaching from above:

\[
\frac{p_{n_i}}{q_{n_i}}>x_B .
\]

For sufficiently large indices:

\[
p_{n_i}
=
\lceil q_{n_i}x_B\rceil .
\]

Therefore:

\[
G_B(q_{n_i})
=
p_{n_i}-q_{n_i}x_B .
\]

Using the continued-fraction bound:

\[
G_B(q_{n_i})
<
\frac{1}{q_{n_i+1}} .
\]

Since:

\[
q_{n_i+1}\rightarrow\infty ,
\]

we obtain:

\[
\lim_{i\rightarrow\infty}
G_B(q_{n_i})
=
0 .
\]

Therefore, the selected convergent subsequence provides block sizes
with vanishing absolute representation gap.

\end{proof}

\begin{remark}

The result is specific to the selected convergent subsequence.

The normalized overhead:

\[
\frac{G_B(K)}{K}
\]

may decrease for many increasing block-size sequences.

However, arbitrary growth of $K$ does not guarantee that the absolute
integer rounding gap decreases.

RICAN provides a deterministic selection rule through continued
fractions that identifies block lengths with exceptionally small
absolute rounding error.

\end{remark}

% =================================================
% V. BLOCK SIZE SELECTION ALGORITHM
% =================================================

\section{Block Size Selection Algorithm}

The theoretical results provide a constructive procedure for selecting
RICAN block sizes.

Given an alphabet cardinality $B$ and a desired representation gap
threshold $\epsilon$, the algorithm computes continued-fraction
convergents of:

\[
\log_2(B).
\]

Each denominator is evaluated as a candidate block size.

For each candidate:

\[
K=q_i,
\]

the required binary length is computed:

\[
n(K)=\lceil K\log_2(B)\rceil .
\]

The representation gap is:

\[
G_B(K)
=
n(K)-K\log_2(B).
\]

For the decimal alphabet:

\[
B=10,
\]

the reference implementation uses:

\[
K=643 .
\]

This block size provides a practical balance between arithmetic
complexity and representation tightness.

Larger convergent-derived values, such as:

\[
K=783032,
\]

demonstrate the asymptotic behavior of the construction but require
larger integer arithmetic operations.

All numerical evaluations of convergents should use arbitrary
precision arithmetic to avoid floating-point rounding errors.

\begin{algorithm}[!t]
\caption{RICAN Block Selection Using Continued Fractions}
\label{alg:block}
\begin{algorithmic}[1]
\REQUIRE Alphabet cardinality $B\geq2$, target gap $\epsilon$
\ENSURE Block size $K$, binary length $n(K)$
\STATE Compute:
\[
x\leftarrow\log_2(B)
\]
\STATE Generate continued-fraction expansion of $x$
\STATE Compute convergents:
\[
p_i/q_i
\]
\FOR{each convergent}
\STATE
\[
K\leftarrow q_i
\]
\STATE
\[
n(K)\leftarrow\lceil Kx\rceil
\]
\STATE
\[
G\leftarrow n(K)-Kx
\]
\IF{$G<\epsilon$}
\RETURN $(K,n(K),G)$
\ENDIF
\ENDFOR
\end{algorithmic}
\end{algorithm}

% =================================================
% VI. REPRESENTATION FRAMEWORK
% =================================================

\section{Representation Framework}

\subsection{General Encoding Model}

The RICAN encoding process constructs an injective mapping from finite
alphabet blocks into binary words.

Given an input sequence:

\[
S=s_0s_1\ldots s_{L-1}
\]

over an alphabet:

\[
\Sigma_B
\]

with cardinality $B$, and a selected block size $K$, the sequence is
partitioned into blocks.

For a complete block:

\[
B_i=(s_0,s_1,\ldots,s_{K-1}),
\]

the corresponding integer value is:

\[
X(B_i)
=
\sum_{j=0}^{K-1}
s_j B^{K-1-j}.
\]

The possible values satisfy:

\[
0\leq X(B_i)<B^K .
\]

Since:

\[
B^K\leq2^{n(K)},
\]

where:

\[
n(K)=\lceil K\log_2(B)\rceil ,
\]

each block can be represented uniquely using exactly $n(K)$ bits.

The final incomplete block, when present, is encoded using its actual
length:

\[
K'<K .
\]

The decoder reconstructs the original sequence by applying the inverse
integer-to-symbol mapping using repeated division by $B$.

\subsection{Injective Representation Mapping}

For a fixed alphabet and block size, the encoder is defined as:

\[
E_B:
\Sigma_B^L
\rightarrow
\{0,1\}^{N(L)} .
\]

The mapping is injective because each valid input block is assigned a
unique integer value.

The decoder is the inverse mapping over the image of the encoder:

\[
D_B:
\operatorname{Im}(E_B)
\rightarrow
\Sigma_B^L .
\]

Therefore:

\[
D_B(E_B(S))=S
\]

for every valid input sequence $S$.

The binary representation space may contain unused states because:

\[
2^{n(K)}>B^K
\]

for non-power-of-two alphabets.

These unused states do not affect reversibility because the decoder
operates only over valid encoded representations.

\subsection{Power-of-Two Alphabets}

For:

\[
B=2^m ,
\]

the representation is exact:

\[
n(K)
=
\lceil Km\rceil
=
Km .
\]

Therefore:

\[
G_B(K)=0 .
\]

Examples include:

\begin{itemize}

\item Binary alphabet ($B=2$): one bit per symbol.

\item DNA alphabet ($B=4$): two bits per symbol.

\item Hexadecimal alphabet ($B=16$): four bits per symbol.

\item Byte alphabet ($B=256$): eight bits per symbol.

\end{itemize}

\subsection{The RJ03 Self-Contained Representation Format}

The reference implementation uses the RJ03 binary representation format.

RJ03 is designed as a self-contained reversible representation
container.

Unlike entropy-coded formats, RJ03 does not require:

\begin{itemize}

\item probability models,

\item adaptive states,

\item external dictionaries,

\item training data.

\end{itemize}

All parameters required for reconstruction are stored inside the
representation file.

The format contains:

\begin{itemize}

\item alphabet definition,

\item alphabet cardinality,

\item block size,

\item original sequence length,

\item encoded block stream.

\end{itemize}

The binary fields use unsigned little-endian representation.

Unused bits in the final byte are padding bits and are ignored during
decoding.

\begin{table}[!t]
\caption{RJ03 Self-Contained Representation Format}
\label{tab:rj03}
\centering
\begin{tabular}{lll}
\toprule
Offset & Size & Field\\
\midrule
0 & 4 & Magic identifier ``RJ03''\\
4 & 4 & Format version\\
8 & 4 & Alphabet cardinality $B$\\
12 & 8 & Block size $K$\\
20 & 8 & Bits per block $n(K)$\\
28 & 4 & Alphabet length\\
32 & variable & Alphabet symbols\\
-- & 8 & Original sequence length $L$\\
-- & 8 & Complete block count\\
-- & 8 & Residual block size\\
-- & variable & Encoded bitstream\\
\bottomrule
\end{tabular}
\end{table}

% =================================================
% VII. EXPERIMENTAL EVALUATION
% =================================================

\section{Experimental Evaluation}

The experiments evaluate the mathematical representation properties of
RICAN.

The objective is not compression benchmarking.

The experiments verify:

\begin{enumerate}

\item exact reconstruction,

\item agreement between theoretical and practical representation sizes,

\item independence of representation length from source statistics,

\item applicability to multiple finite alphabets.

\end{enumerate}

\subsection{Decimal Representation Verification}

A decimal sequence containing:

\[
L=1\,000\,508
\]

digits was encoded with:

\[
B=10
\]

and:

\[
K=643 .
\]

The theoretical block size is:

\[
n(K)=2136
\]

bits.

The measured container size is shown in Table~\ref{tab:decimal-validation}.

\begin{table}[!t]
\caption{Decimal representation verification}
\label{tab:decimal-validation}
\centering
\begin{tabular}{lc}
\toprule
Parameter & Value\\
\midrule
Input symbols & 1,000,508\\
Alphabet size & 10\\
Block size & 643\\
Bits per block & 2136\\
Encoded size & 415,474 bytes\\
Verification & PASS\\
\bottomrule
\end{tabular}
\end{table}

The measured size includes the RJ03 container overhead required for
self-contained reconstruction.

\subsection{Representation-Length Independence}

Because RICAN is a fixed-rate representation framework, the output size
depends on:

\[
B
\]

and:

\[
L .
\]

It does not depend on symbol probability distribution.

Sequences with different statistical properties but identical length
and alphabet cardinality produce identical representation sizes.

\begin{table}[!t]
\caption{Representation-Length Independence}
\label{tab:independence}
\centering
\begin{tabular}{lc}
\toprule
Source Type & Encoded Size\\
\midrule
Mathematical constant digits & Same\\
Random decimal sequence & Same\\
Structured decimal sequence & Same\\
\bottomrule
\end{tabular}
\end{table}

\subsection{Alphabet Generalization}

The same framework applies to multiple finite alphabets.

\begin{table}[!t]
\caption{Finite Alphabet Verification}
\label{tab:alphabet}
\centering
\begin{tabular}{ccc}
\toprule
Alphabet & $B$ & Result\\
\midrule
Binary & 2 & PASS\\
DNA & 4 & PASS\\
Decimal & 10 & PASS\\
Hexadecimal &16& PASS\\
ASCII &128& PASS\\
\bottomrule
\end{tabular}
\end{table}

% =================================================
% VIII. DISCUSSION
% =================================================

\section{Discussion}

\subsection{Main Contribution}

The main contribution of RICAN is the construction of fixed-rate
representation schemes whose integer rounding overhead can be reduced
through continued-fraction block selection.

The existence of fixed-length representations follows directly from
finite-space cardinality.

RICAN contributes a deterministic method for selecting block sizes with
exceptionally small representation gaps.

\subsection{Relationship With Compression}

RICAN does not reduce source information.

It does not exploit:

\begin{itemize}

\item symbol frequencies,

\item repeated patterns,

\item statistical dependencies.

\end{itemize}

Therefore, RICAN should not be considered a replacement for entropy
coding methods.

Its objective is the optimization of deterministic binary
representation overhead.

\subsection{Limitations}

The current framework has several limitations.

Very large convergent block sizes require arbitrary precision integer
arithmetic.

The representation format introduces constant metadata overhead.

Additionally, because the method does not use probability models, it
does not provide compression gains for redundant sources.

% =================================================
% IX. CODE AVAILABILITY
% =================================================

\section{Code Availability and Reproducibility}

The reference implementation, RJ03 specification, and evaluation
scripts are publicly available at:

\begin{center}
\url{https://github.com/SequenceOfManyZerosAndOnes/RICAN}
\end{center}

The repository contains:

\begin{itemize}

\item encoder and decoder implementation,

\item RJ03 format documentation,

\item reproducible examples,

\item mathematical documentation.

\end{itemize}

The experiments were performed using:

\begin{verbatim}
Python 3.13
RICAN RJ03 reference implementation
\end{verbatim}

A permanent archival release through Zenodo is planned.

Software version:

\[
\text{RICAN RJ03 v1.0.0 -- July 2026}
\]

% =================================================
% X. CONCLUSION
% =================================================

\section{Conclusion}

This paper presented RICAN, a fixed-rate binary representation
framework for finite alphabets based on continued-fraction convergent
selection.

The framework applies uniformly to:

\[
\mathcal{A}
=
\{\Sigma_B:B\geq2\}.
\]

For power-of-two alphabets:

\[
B=2^m,
\]

the representation is exact.

For non-power-of-two alphabets, continued-fraction convergents of
$\log_2(B)$ provide block lengths whose absolute representation gap
approaches zero along the selected subsequence.

For decimal representation:

\[
B=10,
\]

the practical implementation uses:

\[
K=643,
\]

while larger convergent-derived values demonstrate the asymptotic
behavior of the construction.

RICAN provides a mathematical framework for studying the boundary
between finite-alphabet representation and binary storage.

Future research directions include:

\begin{itemize}

\item large-scale DNA representation,

\item hardware acceleration,

\item parallel implementations,

\item optimization of very large convergent blocks.

\end{itemize}

% =================================================
% REFERENCES
% =================================================

\bibliographystyle{IEEEtran}
\begin{thebibliography}{00}

\bibitem{shannon1948}
C. E. Shannon,
``A Mathematical Theory of Communication,''
\emph{Bell System Technical Journal},
vol. 27,
pp. 379--423, 623--656,
1948.

\bibitem{huffman1952}
D. A. Huffman,
``A Method for the Construction of Minimum-Redundancy Codes,''
\emph{Proceedings of the IRE},
vol. 40,
no. 9,
pp. 1098--1101,
1952.

\bibitem{ziv1977}
J. Ziv and A. Lempel,
``A Universal Algorithm for Sequential Data Compression,''
\emph{IEEE Transactions on Information Theory},
vol. 23,
no. 3,
pp. 337--343,
1977.

\bibitem{rissanen1979}
J. Rissanen and G. G. Langdon,
``Arithmetic Coding,''
\emph{IBM Journal of Research and Development},
vol. 23,
no. 2,
pp. 149--162,
1979.

\bibitem{duda2013}
J. Duda,
``Asymmetric numeral systems: Entropy coding combining speed of
Huffman coding with compression ratio of arithmetic coding,''
arXiv:1311.2540,
2013.

\bibitem{cover1973}
T. M. Cover,
``Enumerative Source Encoding,''
\emph{IEEE Transactions on Information Theory},
vol. 19,
no. 1,
pp. 73--77,
1973.

\bibitem{immink2004}
K. A. S. Immink,
``A Survey of Enumerative Coding Techniques,''
\emph{IEEE Transactions on Information Theory},
vol. 50,
no. 10,
pp. 2276--2285,
2004.

\bibitem{hardy2008}
G. H. Hardy and E. M. Wright,
\emph{An Introduction to the Theory of Numbers},
6th ed.,
Oxford University Press,
2008.

\bibitem{khinchin1997}
A. Ya. Khinchin,
\emph{Continued Fractions},
Dover Publications,
1997.

\bibitem{cassels1957}
J. W. S. Cassels,
\emph{An Introduction to Diophantine Approximation},
Cambridge University Press,
1957.

\bibitem{rockett1992}
A. M. Rockett and P. Sz\"usz,
\emph{Continued Fractions},
World Scientific,
1992.

\bibitem{brent2010}
R. P. Brent and P. Zimmermann,
\emph{Modern Computer Arithmetic},
Cambridge University Press,
2010.

\end{thebibliography}

\end{document}
